\documentclass[12pt, a4paper]{article}

% ── Packages ──────────────────────────────────────────────────────────────────
\usepackage[utf8]{inputenc}
\usepackage[T1]{fontenc}
\usepackage{amsmath, amssymb, amsthm}
\usepackage{mathtools}
\usepackage{bm}
\usepackage{geometry}
\usepackage{hyperref}
\usepackage{cleveref}
\usepackage{enumitem}
\usepackage{booktabs}
\usepackage{xcolor}
\usepackage{thmtools}
\usepackage{mdframed}

\geometry{margin=2.8cm}
\hypersetup{colorlinks=true, linkcolor=blue!60!black, citecolor=blue!60!black}
\pdfstringdefDisableCommands{%
  \def\ba{a}%
  \def\bx{x}%
  \def\bmu{mu}%
  \def\bm#1{#1}%
}

% ── Theorem environments ───────────────────────────────────────────────────────
\theoremstyle{definition}
\newtheorem{definition}{Definition}[section]
\newtheorem{proposition}{Proposition}[section]
\newtheorem{theorem}{Theorem}[section]
\newtheorem{corollary}{Corollary}[section]
\newtheorem{remark}{Remark}[section]
\newtheorem{lemma}{Lemma}[section]
\renewcommand*{\theHremark}{\thesection.\arabic{remark}}
\renewcommand*{\theHdefinition}{\thesection.\arabic{definition}}
\renewcommand*{\theHproposition}{\thesection.\arabic{proposition}}
\renewcommand*{\theHtheorem}{\thesection.\arabic{theorem}}
\renewcommand*{\theHcorollary}{\thesection.\arabic{corollary}}
\renewcommand*{\theHlemma}{\thesection.\arabic{lemma}}

% ── Custom commands ────────────────────────────────────────────────────────────
\newcommand{\R}{\mathbb{R}}
\newcommand{\E}{\mathbb{E}}
\newcommand{\Var}{\operatorname{Var}}
\newcommand{\Cov}{\operatorname{Cov}}
\newcommand{\N}{\mathcal{N}}
\newcommand{\ba}{\bm{a}}
\newcommand{\bx}{\bm{x}}
\newcommand{\bmu}{\bm{\mu}}
\newcommand{\bz}{\bm{z}}
\newcommand{\Sig}{\Sigma}
\newcommand{\dt}{\Delta_t}
\newcommand{\bc}{\bm{c}}
\newcommand{\bv}{\bm{v}}

% ── Coloured boxes ─────────────────────────────────────────────────────────────
\newmdenv[
  backgroundcolor=blue!5,
  linecolor=blue!40,
  linewidth=0.8pt,
  roundcorner=4pt,
  innertopmargin=8pt,
  innerbottommargin=8pt,
  nobreak=true
]{resultbox}

\newmdenv[
  backgroundcolor=orange!8,
  linecolor=orange!50,
  linewidth=0.8pt,
  roundcorner=4pt,
  innertopmargin=8pt,
  innerbottommargin=8pt,
  nobreak=true
 ]{remarkbox}

\hfuzz=100pt

% ══════════════════════════════════════════════════════════════════════════════
\begin{document}
% ══════════════════════════════════════════════════════════════════════════════

\begin{center}
  {\LARGE\bfseries Joint Score of a Diffusion Model\\[4pt]
  Applied to an AR(1) Sequence}\\[12pt]
  {\large A Self-Contained Derivation}\\[6pt]
  {\normalsize Giovanni Mantovani \& Marco Lomele\quad—\quad
  MSc Thesis, Bocconi University, 2025--2026}\\[4pt]
  {\small Supervisors: Prof.\ Marc M\'ezard \quad Tutor: J\'er\^ome Garnier}
\end{center}

\vspace{0.5cm}
\tableofcontents
\vspace{0.8cm}

% ══════════════════════════════════════════════════════════════════════════════
\section{Motivation and Problem Statement}
% ══════════════════════════════════════════════════════════════════════════════

Standard score-based diffusion models generate \emph{static} objects
(single images). This work studies the \emph{dynamic} setting: a
\textbf{dynamic object} is a sequence of $K+1$ frames
\[
  a_0,\; a_1,\; \ldots,\; a_K \;\in\; \R
\]
indexed by \emph{internal time} $k = 0, 1, \ldots, K$, where consecutive
frames are correlated through a known Markov dynamics.

The naive approach embeds the full trajectory in $\R^{K+1}$ and learns a
single score function — wasteful, because consecutive frames are highly
correlated. Our goal is to understand and exploit the \textbf{temporal
structure of the joint score field}
\[
  \bm{S}(\bx, t)
  \;=\;
  \nabla_{\bx}\log P_t(x_0, x_1, \ldots, x_K),
\]
where $t$ is the diffusion (OU noise) time applied \emph{independently} per
frame, and $\bx = (x_0, \ldots, x_K)$ are the noisy observations.

We pursue this within the \textbf{AR(1) Gaussian toy model}, which is
fully solvable in closed form: every distribution is Gaussian, every
expectation is computable, and the Kalman smoother provides the exact
posterior.

\begin{remarkbox}
\textbf{Critical distinction.} Previous work computed the per-frame
\emph{marginal} score $s(x_u, u, t) = \nabla_{x_u}\log p_t(x_u \mid u)$,
which marginalises over all other frames and therefore discards all
inter-frame correlations. The correct object is the \emph{joint} score
$\bm{S}(\bx, t) \in \R^{K+1}$, whose $k$-th component couples $x_k$ to
all other frames through $\Sigma_t^{-1}$.
\end{remarkbox}

% ══════════════════════════════════════════════════════════════════════════════
\section{The AR(1) Gaussian Model}
% ══════════════════════════════════════════════════════════════════════════════

\subsection{Model specification}

\begin{definition}[AR(1) Gaussian sequence]
The \textbf{clean frames} $\ba = (a_0, a_1, \ldots, a_K) \in \R^{K+1}$
are defined by the following three ingredients, all independent of one
another.
\begin{enumerate}[label=(\roman*)]
  \item \textbf{Initial distribution:}
  \[
    a_0 \;\sim\; \N(\mu_0,\; \sigma_0^2).
  \]
  \item \textbf{Transition rule (AR(1)):} for $k = 0, 1, \ldots, K-1$, where $\alpha \in \R$ is the autoregressive coefficient, where $\alpha \in \R$ is the autoregressive coefficient,
  \[
    a_{k+1} \;=\; \alpha\, a_k \;+\; \eta_k,
    \qquad
    \eta_k \;\overset{\mathrm{i.i.d.}}{\sim}\; \N(0,\; \sigma_\eta^2),
    \qquad
    \eta_k \perp a_0,\; \eta_k \perp \eta_j\; (j \neq k).
  \]
  \item \textbf{Transition kernel:} the above implies
  \[
    M(a_{k+1} \mid a_k) \;=\; \N\!\left(a_{k+1};\; \alpha\,a_k,\;
    \sigma_\eta^2\right).
  \]
\end{enumerate}
\end{definition}

\begin{remark}[Modelling choices]
The zero-mean Gaussian noise $\eta_k$ is the \emph{maximum-entropy}
distribution consistent with a given energy scale $\sigma_\eta^2$,
and it guarantees analytical tractability throughout: sums of Gaussians
are Gaussian, and the entire trajectory $\ba$ remains jointly Gaussian.
The parameter $|\alpha| < 1$ ensures the chain does not explode, with
stationary variance $\sigma_\infty^2 = \sigma_\eta^2/(1-\alpha^2)$.
The i.i.d.\ assumption is on the \emph{noise inputs} $\eta_k$, not on the
frames $a_k$ themselves, which are correlated through the recurrence.
\end{remark}

\subsection{Moments of the chain}

\begin{proposition}[Mean and variance recursions]\label{prop:moments}
The mean $m_k = \E[a_k]$ and variance $\sigma_k^2 = \Var(a_k)$ satisfy
the closed-form recursions:
\[
  m_k \;=\; \alpha^k\,\mu_0,
  \qquad\qquad
  \sigma_k^2 \;=\; \alpha^{2k}\,\sigma_0^2
  \;+\; \sigma_\eta^2\,\frac{1 - \alpha^{2k}}{1 - \alpha^2}.
\]
\end{proposition}

\begin{proof}
\textbf{Mean.} Apply $\E[\cdot]$ to $a_{k+1} = \alpha a_k + \eta_k$.
By linearity and $\E[\eta_k]=0$:
\[
  m_{k+1} = \alpha\,m_k.
\]
Unrolling from $m_0 = \mu_0$ gives $m_k = \alpha^k\mu_0$.

\textbf{Variance.} Since $a_k \perp \eta_k$:
\[
  \sigma_{k+1}^2
  = \Var(\alpha a_k + \eta_k)
  = \alpha^2\,\sigma_k^2 + \sigma_\eta^2.
\]
This is a first-order linear recursion with constant forcing.
Unrolling $k$ times from $\sigma_0^2$:
\[
  \sigma_k^2
  = \alpha^{2k}\sigma_0^2
  + \sigma_\eta^2\sum_{j=0}^{k-1}\alpha^{2j}
  = \alpha^{2k}\sigma_0^2
  + \sigma_\eta^2\,\frac{1-\alpha^{2k}}{1-\alpha^2}. \qquad\square
\]
\end{proof}

\begin{proposition}[Covariance between frames]\label{prop:cov}
For $j \leq k$:
\[
  \Cov(a_j, a_k) \;=\; \alpha^{k-j}\,\sigma_j^2.
\]
Hence the covariance matrix $\Sigma_0 \in \R^{(K+1)\times(K+1)}$ has
entries
\[
  (\Sigma_0)_{jk} \;=\; \alpha^{|j-k|}\,\sigma_{\min(j,k)}^2.
\]
\end{proposition}

\begin{proof}
Unroll the recursion from $a_j$ to $a_k$ ($j \leq k$):
\[
  a_k \;=\; \alpha^{k-j}\,a_j \;+\; \sum_{i=j}^{k-1}\alpha^{k-1-i}\,\eta_i.
\]
The noise terms $\eta_i$ for $i \geq j$ are injected after time $j$, hence
independent of $a_j$. Therefore:
\[
  \Cov(a_j, a_k)
  = \E\!\left[\tilde{a}_j\!\left(\alpha^{k-j}\tilde{a}_j
    + \sum_{i=j}^{k-1}\alpha^{k-1-i}\eta_i\right)\right]
  = \alpha^{k-j}\,\E[\tilde{a}_j^2]
  + \sum_{i=j}^{k-1}\alpha^{k-1-i}\underbrace{\E[\tilde{a}_j\,\eta_i]}_{=\,0}
  = \alpha^{k-j}\,\sigma_j^2,
\]
where $\tilde{a}_j = a_j - m_j$. Symmetry $\Cov(a_j,a_k)=\Cov(a_k,a_j)$
gives the general formula. \qquad$\square$
\end{proof}

\begin{remark}[Non-stationarity]
$\Sigma_0$ depends on position $j$ (not just lag $|j-k|$) because
$\sigma_j^2$ varies with $j$. The chain becomes \emph{stationary}
(i.e.\ $\Cov(a_j,a_k) = \alpha^{|j-k|}\sigma_\infty^2$) if and only if
$\sigma_0^2 = \sigma_\infty^2 = \sigma_\eta^2/(1-\alpha^2)$.
\end{remark}

% ══════════════════════════════════════════════════════════════════════════════
\section{The Clean Joint Distribution \texorpdfstring{$P_0$}{P0}}
% ══════════════════════════════════════════════════════════════════════════════

\subsection{Markov factorisation}

\begin{proposition}[Markov factorisation]\label{prop:markov}
The joint density of the clean trajectory factorises as
\[
  P_0(a_0, a_1, \ldots, a_K)
  \;=\;
  p_0(a_0)\;\prod_{k=0}^{K-1} M(a_{k+1} \mid a_k),
\]
or explicitly:
\[
  P_0(\ba)
  \;=\;
  \frac{1}{(2\pi)^{(K+1)/2}\,\sigma_0\,\sigma_\eta^K}
  \exp\!\left(
    -\frac{(a_0-\mu_0)^2}{2\sigma_0^2}
    \;-\;
    \sum_{k=0}^{K-1}\frac{(a_{k+1}-\alpha\,a_k)^2}{2\sigma_\eta^2}
  \right).
\]
\end{proposition}

\begin{proof}
The chain rule of probability gives:
\[
  P_0(\ba) = p(a_0)\cdot p(a_1\mid a_0)\cdot p(a_2\mid a_1,a_0)
  \cdots p(a_K\mid a_{K-1},\ldots,a_0).
\]
The Markov property reduces each factor: $p(a_k\mid a_{k-1},\ldots,a_0)
= M(a_k\mid a_{k-1})$. Substituting the Gaussian densities gives the
explicit product form. \qquad$\square$
\end{proof}

\subsection{The precision matrix}

The exponent of $P_0(\ba)$ is quadratic in $\ba$. Define the
\emph{precision matrix} $\Sigma_0^{-1}$ via
\[
  \frac{(a_0-\mu_0)^2}{2\sigma_0^2}
  + \sum_{k=0}^{K-1}\frac{(a_{k+1}-\alpha\,a_k)^2}{2\sigma_\eta^2}
  \;\equiv\;
  \frac{1}{2}(\ba-\bmu_a)^\top \Sigma_0^{-1}(\ba-\bmu_a),
\]
where $\bmu_a = (\mu_0, \alpha\mu_0, \ldots, \alpha^K\mu_0)^\top$.

\begin{proposition}[Tridiagonal precision matrix]\label{prop:tridiag}
The precision matrix $\Sigma_0^{-1}$ is \textbf{symmetric} and
\textbf{tridiagonal} with entries:
\begin{align*}
  (\Sigma_0^{-1})_{00}
    &= \frac{1}{\sigma_0^2} + \frac{\alpha^2}{\sigma_\eta^2}, \\
  (\Sigma_0^{-1})_{kk}
    &= \frac{1+\alpha^2}{\sigma_\eta^2},
    \quad k = 1, \ldots, K-1, \\
  (\Sigma_0^{-1})_{KK}
    &= \frac{1}{\sigma_\eta^2}, \\
  (\Sigma_0^{-1})_{k,k+1}
  = (\Sigma_0^{-1})_{k+1,k}
    &= -\frac{\alpha}{\sigma_\eta^2},
    \quad k = 0, \ldots, K-1, \\
  (\Sigma_0^{-1})_{jk}
    &= 0,
    \quad |j-k| \geq 2.
\end{align*}
\end{proposition}

\begin{proof}
Expand the exponent and collect terms. For a symmetric matrix $A$, the
identity
\[
  \mathbf{v}^\top A \mathbf{v}
  = \sum_i A_{ii}v_i^2 + 2\sum_{i<j}A_{ij}v_iv_j
\]
implies: $A_{ii}$ equals the coefficient of $v_i^2$, and $A_{ij}$ ($i<j$)
equals \emph{half} the coefficient of $v_iv_j$.

Setting $\tilde{a}_k = a_k - \alpha^k\mu_0$ (centred variables):

\begin{itemize}
  \item \textbf{Coefficient of $\tilde{a}_0^2$:} appears in
    $(a_0-\mu_0)^2/\sigma_0^2$ as $1/\sigma_0^2$, and in
    $(a_1 - \alpha a_0)^2/\sigma_\eta^2$ as $\alpha^2/\sigma_\eta^2$.
    Sum: $(\Sigma_0^{-1})_{00} = 1/\sigma_0^2 + \alpha^2/\sigma_\eta^2$.

  \item \textbf{Coefficient of $\tilde{a}_k^2$, interior $k$:} frame $k$
    appears as $\tilde{a}_k^2$ in transition $k-1$ and as
    $\alpha^2\tilde{a}_k^2$ in transition $k$. Sum:
    $(\Sigma_0^{-1})_{kk} = (1+\alpha^2)/\sigma_\eta^2$.

  \item \textbf{Coefficient of $\tilde{a}_K^2$:} only in the last
    transition as $\tilde{a}_K^2/\sigma_\eta^2$.
    Hence $(\Sigma_0^{-1})_{KK} = 1/\sigma_\eta^2$.

  \item \textbf{Coefficient of $\tilde{a}_k\tilde{a}_{k+1}$:} the cross
    term $-2\alpha\tilde{a}_{k+1}\tilde{a}_k/\sigma_\eta^2$ in transition
    $k$ gives, after halving, $(\Sigma_0^{-1})_{k,k+1} = -\alpha/\sigma_\eta^2$.

  \item \textbf{Terms $\tilde{a}_j\tilde{a}_k$ with $|j-k|\geq 2$:}
    such cross-products never appear in the expanded sum of squared
    residuals — a direct algebraic consequence of the Markov property.
    Hence $(\Sigma_0^{-1})_{jk} = 0$. \qquad$\square$
\end{itemize}
\end{proof}

\begin{remark}[Conditional independence interpretation]
For a jointly Gaussian vector, $(\Sigma_0^{-1})_{jk} = 0$ if and only if
$a_j \perp a_k \mid \{a_\ell : \ell \neq j, k\}$. The tridiagonal
structure of $\Sigma_0^{-1}$ is therefore the algebraic encoding of the
Markov property: conditioning on all intermediate frames screens off
$a_j$ and $a_k$ for $|j-k|\geq 2$. The covariance matrix $\Sigma_0$ is
\emph{dense} (all entries nonzero for $\alpha\neq 0$), encoding
\emph{total} correlations including all indirect paths through
intermediaries; the precision matrix encodes only \emph{direct} couplings.
More precisely, from the conditional distribution $a_j \mid \{a_k\}_{k\neq j}$,
the conditional mean satisfies
\[
  \E\!\left[a_j \mid \{a_k\}_{k\neq j}\right]
  = \mu_j - \frac{1}{(\Sigma_0^{-1})_{jj}}\sum_{k\neq j}
    (\Sigma_0^{-1})_{jk}(a_k - \mu_k),
\]
so a zero off-diagonal entry $(\Sigma_0^{-1})_{jk}=0$ means $a_k$ does not
appear in the conditional mean of $a_j$ given everything else.
\end{remark}

\begin{resultbox}
\textbf{Clean joint distribution.} The AR(1) trajectory is a
$(K+1)$-dimensional Gaussian:
\[
  \boxed{
  P_0(\ba)
  = \frac{1}{(2\pi)^{(K+1)/2}\,|\Sigma_0|^{1/2}}
    \exp\!\left(
      -\frac{1}{2}(\ba - \bmu_a)^\top\,\Sigma_0^{-1}\,(\ba - \bmu_a)
    \right),}
\]
with mean $\bmu_a = (\mu_0, \alpha\mu_0, \ldots, \alpha^K\mu_0)^\top$
and tridiagonal precision matrix $\Sigma_0^{-1}$ given in
Proposition~\ref{prop:tridiag}.
\end{resultbox}

% ══════════════════════════════════════════════════════════════════════════════
\section{The Noisy Joint Distribution \texorpdfstring{$P_t$}{Pt}}
% ══════════════════════════════════════════════════════════════════════════════

\subsection{The Ornstein-Uhlenbeck forward process}

\begin{definition}[OU diffusion per frame]
At diffusion time $t \geq 0$, each clean frame $a_k$ is independently
corrupted:
\[
  x_k \;=\; a_k\,e^{-t} \;+\; \sqrt{\dt}\;z_k,
  \qquad
  z_k \overset{\mathrm{i.i.d.}}{\sim} \N(0,1),
  \quad
  \bm{z} \perp \ba,
  \qquad
  \dt \;=\; 1 - e^{-2t}.
\]
Note the identity $e^{-2t} + \dt = 1$: signal energy plus noise energy
is conserved.
\end{definition}

The scaling rule $\sqrt{\dt}\,z_k \sim \N(0, \dt)$ gives:
\[
  x_k \mid a_k \;\sim\; \N\!\left(a_k e^{-t},\; \dt\right),
  \qquad
  p(x_k\mid a_k)
  = \frac{1}{\sqrt{2\pi\dt}}
    \exp\!\left(-\frac{(x_k - a_k e^{-t})^2}{2\dt}\right).
\]

Since the noises $z_k$ are independent across frames, the joint
conditional likelihood factorises:
\[
  p(\bx \mid \ba)
  \;=\;
  \prod_{k=0}^{K}\,\N\!\left(x_k;\; a_k e^{-t},\; \dt\right).
\]

\subsection{The marginal distribution of noisy observations}

The noisy observations $\bx = (x_0, \ldots, x_K)$ are obtained by
marginalising over the hidden clean frames:

\begin{resultbox}
\[
  P_t(\bx)
  \;=\;
  \int_{\R^{K+1}} P_0(\ba)\;\prod_{k=0}^{K}
  \N\!\left(x_k;\; a_k e^{-t},\; \dt\right)
  \;d\ba.
  \tag{Mézard's board formula}
\]
\end{resultbox}

We compute this integral by two independent methods.

\subsection{Method 1: Linear transformation of Gaussians}

\begin{theorem}[$P_t$ is Gaussian]\label{thm:Pt_linear}
The marginal $P_t(\bx)$ is a $(K+1)$-dimensional Gaussian:
\[
  P_t(\bx) \;=\; \N\!\left(\bx;\; \bmu_t,\; \Sigma_t\right),
\]
with
\[
  \bmu_t \;=\; e^{-t}\,\bmu_a,
  \qquad
  \Sigma_t \;=\; e^{-2t}\,\Sigma_0 \;+\; \dt\, I_{K+1}.
\]
\end{theorem}

\begin{proof}
Write the observation model in vector form:
\[
  \bx = e^{-t}\,\ba + \sqrt{\dt}\,\bz,
  \qquad
  \ba \sim \N(\bmu_a, \Sigma_0),
  \quad
  \bz \sim \N(\bm{0}, I_{K+1}),
  \quad
  \ba \perp \bz.
\]
$\bx$ is a linear function of two independent Gaussian vectors,
hence Gaussian.

\textbf{Mean:}
\[
  \bmu_t
  = \E[\bx]
  = e^{-t}\,\E[\ba] + \sqrt{\dt}\,\E[\bz]
  = e^{-t}\,\bmu_a.
\]

\textbf{Covariance:} using independence $\ba \perp \bz$ and the
bilinearity of covariance,
\[
  \Sigma_t
  = \Cov\!\left(e^{-t}\,\ba + \sqrt{\dt}\,\bz\right)
  = e^{-2t}\,\Cov(\ba) + \dt\,\Cov(\bz)
  = e^{-2t}\,\Sigma_0 + \dt\,I_{K+1}. \qquad\square
\]
\end{proof}

\begin{remark}[Physical interpretation of $\Sigma_t$]
The covariance $\Sigma_t = e^{-2t}\Sigma_0 + \dt I$ has two competing
contributions:
\begin{itemize}
  \item $e^{-2t}\Sigma_0$: the \emph{signal} covariance, carrying all
    inter-frame AR(1) correlations, shrinking exponentially with $t$.
  \item $\dt I$: the \emph{noise} covariance, diagonal (independent
    noise per frame), growing from $0$ to $1$ as $t$ increases.
\end{itemize}
At $t=0$: $\Sigma_0 = \Sigma_0$ (pure signal). As $t\to\infty$:
$\Sigma_t \to I$ (pure noise, all correlation erased).
\end{remark}

\subsection{Method 2: Direct Gaussian integral}

We now verify Theorem~\ref{thm:Pt_linear} by explicitly computing
the integral, which is instructive in its own right.

\begin{proof}[Proof of Theorem~\ref{thm:Pt_linear} by direct integration]
We compute $P_t(\bx) = \int_{\R^{K+1}} P_0(\ba)\,p(\bx\mid\ba)\,d\ba$
by treating the integrand as a Gaussian in $\ba$ and completing the square.

\medskip
\noindent\textit{Step 1: combine the two exponentials.}
Up to positive normalisation constants (absorbed into the proportionality),
\[
  P_0(\ba)\cdot p(\bx\mid\ba)
  \;\propto\;
  \exp\!\left(
    -\frac{1}{2}(\ba-\bmu_a)^\top\Sigma_0^{-1}(\ba-\bmu_a)
    -\frac{1}{2\dt}\|\bx - e^{-t}\ba\|^2
  \right).
\]
Expand the second quadratic (writing $\|\bx - e^{-t}\ba\|^2 =
\bx^\top\bx - 2e^{-t}\bx^\top\ba + e^{-2t}\ba^\top\ba$):
\[
  \frac{1}{2\dt}\|\bx - e^{-t}\ba\|^2
  = \frac{e^{-2t}}{2\dt}\,\ba^\top\ba
    - \frac{e^{-t}}{\dt}\,\bx^\top\ba
    + \frac{1}{2\dt}\,\bx^\top\bx.
\]
The last term does not depend on $\ba$; absorb it into the proportionality.
Expand the first quadratic similarly:
\[
  \frac{1}{2}(\ba-\bmu_a)^\top\Sigma_0^{-1}(\ba-\bmu_a)
  = \frac{1}{2}\ba^\top\Sigma_0^{-1}\ba
    - \ba^\top\Sigma_0^{-1}\bmu_a
    + \frac{1}{2}\bmu_a^\top\Sigma_0^{-1}\bmu_a.
\]
The last term is a constant in $\ba$; absorb it. Collecting all $\ba$-dependent
terms in the exponent:
\[
  -\frac{1}{2}\,\ba^\top
    \underbrace{\left(\Sigma_0^{-1} + \frac{e^{-2t}}{\dt}\,I\right)}_{\Lambda}\,\ba
  + \ba^\top
    \underbrace{\left(\Sigma_0^{-1}\bmu_a + \frac{e^{-t}}{\dt}\,\bx\right)}_{\bm{h}}.
\]

\medskip
\noindent\textit{Step 2: complete the square in $\ba$.}
Define
\begin{align*}
  \Lambda &\;=\; \Sigma_0^{-1} + \frac{e^{-2t}}{\dt}\,I
  \quad\in\R^{(K+1)\times(K+1)},\\
  \bm{h} &\;=\; \Sigma_0^{-1}\bmu_a + \frac{e^{-t}}{\dt}\,\bx,\\
  \bm{m} &\;=\; \Lambda^{-1}\bm{h}
  \quad\text{(posterior mean of $\ba$ given $\bx$)}.
\end{align*}
The standard identity $-\frac{1}{2}\ba^\top\Lambda\ba + \ba^\top\bm{h}
= -\frac{1}{2}(\ba-\bm{m})^\top\Lambda(\ba-\bm{m})
+ \frac{1}{2}\bm{m}^\top\Lambda\bm{m}$
rewrites the exponent as a Gaussian in $\ba$ centred at $\bm{m}$,
plus a term $\frac{1}{2}\bm{m}^\top\Lambda\bm{m}$ that depends only on $\bx$.

\medskip
\noindent\textit{Step 3: perform the Gaussian integral over $\ba$.}
\[
  \int_{\R^{K+1}}
    \exp\!\left(-\tfrac{1}{2}(\ba-\bm{m})^\top\Lambda(\ba-\bm{m})\right)d\ba
  = (2\pi)^{(K+1)/2}\,|\Lambda|^{-1/2}.
\]
The remaining $\bx$-dependent part of the exponent is:
\[
  Q(\bx)
  \;\equiv\;
  \frac{1}{2}\bm{m}^\top\Lambda\bm{m} - \frac{1}{2\dt}\,\bx^\top\bx
  = \frac{1}{2}\bm{h}^\top\Lambda^{-1}\bm{h} - \frac{1}{2\dt}\,\bx^\top\bx.
\]

\medskip
\noindent\textit{Step 4: identify $Q(\bx)$ as a Gaussian quadratic in $\bx$.}
Substitute $\bm{h} = \Sigma_0^{-1}\bmu_a + \frac{e^{-t}}{\dt}\bx$ and
apply the \textbf{Woodbury identity}
$\Lambda^{-1} = \left(\Sigma_0^{-1}+\frac{e^{-2t}}{\dt}I\right)^{-1}
= \Sigma_0 - e^{-2t}\Sigma_0\Sigma_t^{-1}\Sigma_0$
(derivable from $(A+UCV)^{-1} = A^{-1}-A^{-1}U(C^{-1}+VA^{-1}U)^{-1}VA^{-1}$
with $A=\Sigma_0^{-1}$, $U=V^\top=I$, $C=\frac{e^{-2t}}{\dt}I$).
After expanding and collecting the $\bx^\top(\cdot)\bx$ terms, one finds
that all cross-terms reduce to:
\[
  Q(\bx)
  = -\frac{1}{2}(\bx - e^{-t}\bmu_a)^\top \Sigma_t^{-1}(\bx - e^{-t}\bmu_a)
    + \text{const},
\]
where $\Sigma_t = e^{-2t}\Sigma_0 + \dt I$ (one verifies:
$\frac{1}{\dt}I - \frac{e^{-2t}}{\dt^2}\Sigma_0\Lambda^{-1} = \Sigma_t^{-1}$
by multiplying both sides by $\Sigma_t$ and using $\Lambda = \Sigma_0^{-1}+\frac{e^{-2t}}{\dt}I$).

\medskip
\noindent\textit{Conclusion.}
Assembling all factors, $P_t(\bx) \propto
\exp\!\left(-\frac{1}{2}(\bx-\bmu_t)^\top\Sigma_t^{-1}(\bx-\bmu_t)\right)$
with $\bmu_t = e^{-t}\bmu_a$ and $\Sigma_t = e^{-2t}\Sigma_0+\dt I$,
confirming it is $\N(\bmu_t,\Sigma_t)$.
\end{proof}

\begin{resultbox}
\textbf{Noisy joint distribution.}
\[
  \boxed{
  P_t(\bx)
  = \frac{1}{(2\pi)^{(K+1)/2}|\Sigma_t|^{1/2}}
    \exp\!\left(
      -\frac{1}{2}(\bx - \bmu_t)^\top\,\Sigma_t^{-1}\,(\bx - \bmu_t)
    \right),}
\]
with $\bmu_t = e^{-t}\bmu_a$ and
$\Sigma_t = e^{-2t}\Sigma_0 + \dt I_{K+1}$.
\end{resultbox}

% ══════════════════════════════════════════════════════════════════════════════
\section{The Joint Score Function}
% ══════════════════════════════════════════════════════════════════════════════

\subsection{Derivation}

\begin{definition}[Joint score]
The \textbf{joint score} at diffusion time $t$ is:
\[
  \bm{S}(\bx, t)
  \;=\;
  \nabla_{\bx}\log P_t(\bx)
  \;\in\; \R^{K+1}.
\]
Its $k$-th component $S_k(\bx,t) = \partial_{x_k}\log P_t(\bx)$
measures how increasing $x_k$ changes the log-probability of the
entire noisy trajectory.
\end{definition}

\begin{theorem}[Joint score is linear]\label{thm:score}
\[
  \bm{S}(\bx, t)
  \;=\;
  -\Sigma_t^{-1}\,\bigl(\bx - \bmu_t\bigr).
\]
\end{theorem}

\begin{proof}
Since $P_t(\bx) = \N(\bx;\bmu_t,\Sigma_t)$, its log-density is:
\[
  \log P_t(\bx)
  = -\frac{1}{2}(\bx-\bmu_t)^\top\Sigma_t^{-1}(\bx-\bmu_t)
    \underbrace{- \frac{K+1}{2}\log(2\pi) - \frac{1}{2}\log|\Sigma_t|}_{\text{constant in }\bx}.
\]
It therefore suffices to differentiate $-\frac{1}{2}(\bx-\bmu_t)^\top\Sigma_t^{-1}(\bx-\bmu_t)$
with respect to $\bx$. Expand the quadratic form using bilinearity:
\[
  (\bx-\bmu_t)^\top\Sigma_t^{-1}(\bx-\bmu_t)
  = \bx^\top\Sigma_t^{-1}\bx
    - 2\,\bmu_t^\top\Sigma_t^{-1}\bx
    + \underbrace{\bmu_t^\top\Sigma_t^{-1}\bmu_t}_{\text{constant in }\bx}.
\]
We differentiate the two remaining $\bx$-dependent terms separately.

\medskip
\noindent\textit{Step 1: gradient of the quadratic term.}
For any symmetric matrix $A \in \R^{n\times n}$ and $\bm{v}\in\R^n$,
compute the $m$-th partial derivative of $\bm{v}^\top A\bm{v}
= \sum_{i,j}A_{ij}v_iv_j$ by the product rule:
\[
  \frac{\partial}{\partial v_m}\Bigl(\bm{v}^\top A\bm{v}\Bigr)
  = \sum_j A_{mj}v_j \;+\; \sum_i A_{im}v_i
  = \bigl(A\bm{v}\bigr)_m + \bigl(A^\top\bm{v}\bigr)_m.
\]
Since $\Sigma_t^{-1}$ is symmetric ($\Sigma_t^{-1} = (\Sigma_t^{-1})^\top$,
as the inverse of a symmetric matrix is symmetric), the two sums are
equal and:
\[
  \nabla_{\bx}\!\left(\bx^\top\Sigma_t^{-1}\bx\right)
  = 2\,\Sigma_t^{-1}\bx.
\]

\medskip
\noindent\textit{Step 2: gradient of the linear term.}
The term $-2\,\bmu_t^\top\Sigma_t^{-1}\bx$ is linear in $\bx$.
Write it as $\bm{c}^\top\bx$ where $\bm{c} = -2\,\Sigma_t^{-1}\bmu_t$
is a constant vector (independent of $\bx$). Since
$\nabla_{\bx}(\bm{c}^\top\bx) = \bm{c}$ for any constant $\bm{c}$:
\[
  \nabla_{\bx}\!\left(-2\,\bmu_t^\top\Sigma_t^{-1}\bx\right)
  = -2\,\Sigma_t^{-1}\bmu_t.
\]

\medskip
\noindent\textit{Combining Steps 1 and 2:}
\begin{align*}
  \bm{S}(\bx,t)
  &= \nabla_{\bx}\!\left[\log P_t(\bx)\right] \\
  &= -\frac{1}{2}\,\nabla_{\bx}\!\left(
      \bx^{\top}\Sigma_t^{-1}\bx \;-\; 2\,\bmu_t^{\top}\Sigma_t^{-1}\bx
    \right) \\
  &= -\frac{1}{2}\!\left(2\,\Sigma_t^{-1}\bx - 2\,\Sigma_t^{-1}\bmu_t\right) \\
  &= -\Sigma_t^{-1}(\bx - \bmu_t).
\end{align*}
\end{proof}

\begin{remark}[Linearity is a Gaussian property]
The score is affine in $\bx$:
\[
  \bm{S}(\bx, t) = \underbrace{-\Sigma_t^{-1}}_{A(t)}\,\bx
  + \underbrace{\Sigma_t^{-1}\bmu_t}_{\bm{b}(t)}.
\]
This linearity is specific to Gaussian distributions. For a non-Gaussian
prior $P_0$ (e.g.\ a mixture of Gaussians), $P_t$ is non-Gaussian for
small $t$ and the score is a nonlinear function of $\bx$. The AR(1)
Gaussian model is the unique tractable case where linearity holds exactly
at all $t$.
\end{remark}

\subsection{The \texorpdfstring{$k$}{k}-th component and inter-frame coupling}

The $k$-th component of the joint score is:

\begin{corollary}[Component-wise joint score]
\[
  S_k(\bx, t)
  \;=\;
  -\sum_{j=0}^{K}(\Sigma_t^{-1})_{kj}\,(x_j - \mu_{t,j}),
\]
where $\mu_{t,j} = e^{-t}\alpha^j\mu_0$. This couples $x_k$ to
\textbf{all other frames} $x_j$ through the precision matrix entries
$(\Sigma_t^{-1})_{kj}$.
\end{corollary}

\begin{remark}[Limiting regimes]
\hfill
\begin{itemize}
  \item \textbf{$t \to 0$ (clean signal):} $\Sigma_t^{-1} \to
    \Sigma_0^{-1}$, which is tridiagonal. Hence
    $S_k$ depends only on $x_{k-1}$, $x_k$, $x_{k+1}$ —
    nearest-neighbour coupling only.

  \item \textbf{$t \to \infty$ (pure noise):} $\Sigma_t^{-1}
    \to \dt^{-1}I \to I$, so $S_k \approx -x_k/\dt$ — each
    frame's score depends only on itself, all inter-frame coupling
    vanishes.

  \item \textbf{Intermediate $t$:} $\Sigma_t^{-1}$ is not tridiagonal
    in general. Adding $\dt I$ to $\Sigma_0$ and inverting fills in
    the structural zeros of $\Sigma_0^{-1}$, producing small but nonzero
    off-tridiagonal entries. The tridiagonal structure is only approximate
    for $t > 0$.
\end{itemize}
\end{remark}

% ══════════════════════════════════════════════════════════════════════════════
\section{Connection to M\'ezard's Formula}
% ══════════════════════════════════════════════════════════════════════════════

The formula derived in Theorem~\ref{thm:score} has an equivalent
probabilistic representation as a conditional expectation, which is what
M\'ezard wrote on the board during the meeting.

\begin{theorem}[Score as a conditional expectation]\label{thm:mezard}
The $k$-th component of the joint score satisfies:
\[
  \boxed{
  S_k(\bx, t)
  \;=\;
  \left\langle
    \frac{a_k\,e^{-t} - x_k}{\dt}
  \right\rangle_{\!\ba \mid \bx},}
\]
where $\langle \cdot \rangle_{\ba \mid \bx}$ denotes expectation under
the posterior
$P(\ba \mid \bx) \propto P_0(\ba)\prod_{j}\N(x_j;\,a_j e^{-t},\,\dt)$.
\end{theorem}

\begin{proof}
By definition:
\[
  S_k(\bx,t)
  = \partial_{x_k}\log P_t(\bx)
  = \frac{\partial_{x_k} P_t(\bx)}{P_t(\bx)}.
\]
Write $P_t(\bx) = \int P_0(\ba)\prod_j \N(x_j;\,a_j e^{-t},\,\dt)\,d\ba$.
Only the $k$-th factor depends on $x_k$. Using the identity
$\partial_{x}\log\N(x;\,\mu,\,\sigma^2) = (\mu-x)/\sigma^2$:
\[
  \partial_{x_k}\N(x_k;\,a_k e^{-t},\,\dt)
  = \N(x_k;\,a_k e^{-t},\,\dt)\cdot\frac{a_k e^{-t} - x_k}{\dt}.
\]
Therefore:
\[
  \partial_{x_k}P_t(\bx)
  = \int P_0(\ba)\prod_j\N(x_j;\,a_j e^{-t},\,\dt)
    \cdot\frac{a_k e^{-t}-x_k}{\dt}\,d\ba.
\]
Dividing by $P_t(\bx)$ turns the integral into an expectation under
the posterior $P(\ba\mid\bx)$:
\[
  S_k(\bx,t)
  = \left\langle\frac{a_k e^{-t} - x_k}{\dt}\right\rangle_{\ba\mid\bx}.
  \qquad\square
\]
\end{proof}

\begin{remark}[Denoising interpretation]
The formula $S_k = (e^{-t}\langle a_k\rangle_{\ba\mid\bx} - x_k)/\dt$
is a \textbf{denoising residual}: the expected clean signal scaled back
to the observed level, minus the actual noisy observation, divided by
the noise variance. Three features deserve emphasis:
\begin{enumerate}
  \item The expectation $\langle a_k\rangle_{\ba\mid\bx}$ conditions on
    the \textbf{full noisy trajectory} $\bx$, not just $x_k$. This is
    the key difference from the marginal score.
  \item For Gaussian AR(1), the posterior $P(\ba\mid\bx)$ is Gaussian,
    and the posterior mean $\langle\ba\rangle_{\ba\mid\bx}$ is
    computable exactly by the \textbf{Kalman smoother}.
  \item The posterior mean is \textbf{linear} in $\bx$:
    $\langle\ba\rangle_{\ba\mid\bx} = e^{-t}\Sigma_0\Sigma_t^{-1}
    (\bx - \bmu_t) + \bmu_a$, which is consistent with the
    explicit formula $S_k = -[\Sigma_t^{-1}(\bx-\bmu_t)]_k$.
\end{enumerate}
\end{remark}

\subsection{Equivalence of the two formulas}

We now prove algebraically that M\'ezard's conditional expectation formula
and the explicit linear formula of Theorem~\ref{thm:score} are identical.
The key ingredient is the \textbf{posterior mean} of the clean frames.

\begin{lemma}[Gaussian posterior mean]\label{lem:postmean}
For Gaussian $P_0(\ba) = \N(\ba;\bmu_a,\Sigma_0)$ and Gaussian
likelihood $p(\bx\mid\ba) = \prod_k\N(x_k;\,a_ke^{-t},\dt)$,
the posterior $P(\ba\mid\bx)$ is Gaussian with mean:
\[
  \langle\ba\rangle_{\ba\mid\bx}
  \;=\;
  \bmu_a \;+\; e^{-t}\,\Sigma_0\,\Sigma_t^{-1}\,(\bx - \bmu_t).
\]
\end{lemma}

\begin{proof}
The posterior precision is $\Lambda = \Sigma_0^{-1}+\frac{e^{-2t}}{\dt}I$
(from Step~2 of the direct integral proof) and the posterior mean is
$\bm{m} = \Lambda^{-1}(\Sigma_0^{-1}\bmu_a + \frac{e^{-t}}{\dt}\bx)$.
Apply the Woodbury identity $\Lambda^{-1} = \Sigma_0 - e^{-2t}\Sigma_0\Sigma_t^{-1}\Sigma_0$
and the relation $\frac{e^{-2t}}{\dt}\Sigma_0 = \Sigma_t^{-1}\Sigma_0\cdot\frac{e^{-2t}\Sigma_0}{\dt}$...

More directly: use the standard Gaussian conditioning formula.
Because $\bx = e^{-t}\ba + \sqrt{\dt}\bz$ with $\ba\perp\bz$,
the joint covariance between $\ba$ and $\bx$ is
$\Cov(\ba,\bx) = e^{-t}\Cov(\ba,\ba) = e^{-t}\Sigma_0$.
The conditional mean formula for jointly Gaussian variables gives:
\[
  \langle\ba\rangle_{\ba\mid\bx}
  = \bmu_a + \Cov(\ba,\bx)\,[\Cov(\bx,\bx)]^{-1}(\bx-\bmu_t)
  = \bmu_a + e^{-t}\Sigma_0\,\Sigma_t^{-1}(\bx-\bmu_t). \qquad\square
\]
\end{proof}

\begin{proposition}[Equivalence of the two score formulas]
The M\'ezard formula and the linear formula agree exactly:
\[
  \left\langle\frac{a_k e^{-t}-x_k}{\dt}\right\rangle_{\!\ba\mid\bx}
  \;=\;
  -\bigl[\Sigma_t^{-1}(\bx-\bmu_t)\bigr]_k.
\]
\end{proposition}

\begin{proof}
Since $x_k$ is a constant under $\langle\cdot\rangle_{\ba\mid\bx}$:
\[
  \left\langle\frac{a_k e^{-t}-x_k}{\dt}\right\rangle_{\!\ba\mid\bx}
  = \frac{e^{-t}\langle a_k\rangle_{\ba\mid\bx} - x_k}{\dt}.
\]
Substitute Lemma~\ref{lem:postmean} for the $k$-th component:
\[
  = \frac{e^{-t}\!\left(\mu_{a,k} + e^{-t}[\Sigma_0\Sigma_t^{-1}(\bx-\bmu_t)]_k\right) - x_k}{\dt}.
\]
Use $e^{-t}\mu_{a,k} = \mu_{t,k}$ (since $\bmu_t = e^{-t}\bmu_a$) and write
$e^{-2t}\Sigma_0 = \Sigma_t - \dt I$ (rearranging $\Sigma_t = e^{-2t}\Sigma_0+\dt I$),
so $e^{-2t}\Sigma_0\Sigma_t^{-1} = (\Sigma_t-\dt I)\Sigma_t^{-1} = I - \dt\Sigma_t^{-1}$:
\[
  = \frac{\mu_{t,k} + [I - \dt\Sigma_t^{-1}]_k(\bx-\bmu_t) - x_k}{\dt}.
\]
Expanding $[I - \dt\Sigma_t^{-1}]_k(\bx-\bmu_t)
= (x_k - \mu_{t,k}) - \dt[\Sigma_t^{-1}(\bx-\bmu_t)]_k$:
\[
  = \frac{\mu_{t,k} + (x_k - \mu_{t,k}) - \dt[\Sigma_t^{-1}(\bx-\bmu_t)]_k - x_k}{\dt}
  = \frac{-\dt[\Sigma_t^{-1}(\bx-\bmu_t)]_k}{\dt}
  = -[\Sigma_t^{-1}(\bx-\bmu_t)]_k. \qquad\square
\]
\end{proof}

% ══════════════════════════════════════════════════════════════════════════════
\section{Summary and Notation Table}
% ══════════════════════════════════════════════════════════════════════════════

\subsection{The complete derivation chain}

\[
  \underbrace{P_0(\ba)}_{\text{AR(1) prior}}
  \;\xrightarrow{\;\text{OU noise, independent per frame}\;}
  \underbrace{P_t(\bx) = \N(\bmu_t,\,\Sigma_t)}_{\text{noisy joint dist.}}
  \;\xrightarrow{\;\nabla_{\bx} \log\;}
  \underbrace{\bm{S}(\bx,t) = -\Sigma_t^{-1}(\bx-\bmu_t)}_{\text{joint score}}
\]

\subsection{Key formulas collected}

\begin{center}
\renewcommand{\arraystretch}{1.6}
\begin{tabular}{lll}
\toprule
\textbf{Object} & \textbf{Formula} & \textbf{Key property} \\
\midrule
Markov prior & $P_0(\ba) = p_0(a_0)\prod_k M(a_{k+1}\mid a_k)$ &
  Markov factorisation \\
Covariance & $(\Sigma_0)_{jk} = \alpha^{|j-k|}\sigma_{\min(j,k)}^2$ &
  Dense \\
Precision & $(\Sigma_0^{-1})_{k,k\pm1} = -\alpha/\sigma_\eta^2$ &
  Tridiagonal (sparse) \\
OU variance & $\dt = 1 - e^{-2t}$ & $e^{-2t}+\dt=1$ \\
Noisy covariance & $\Sigma_t = e^{-2t}\Sigma_0 + \dt I$ &
  Signal + noise \\
Noisy mean & $\bmu_t = e^{-t}\bmu_a$ & Shrinks by $e^{-t}$ \\
Joint score (algebraic) & $\bm{S} = -\Sigma_t^{-1}(\bx-\bmu_t)$ &
  Linear in $\bx$ \\
Joint score (M\'ezard) & $S_k = \langle(a_k e^{-t}-x_k)/\dt\rangle_{\ba\mid\bx}$ &
  Denoising residual \\
Marginal score (old) & $s(x_k,k,t) = -(x_k-\mu_{t,k})/T_{kk}$ &
  Ignores other frames \\
\bottomrule
\end{tabular}
\end{center}

\subsection{Notation reference}

\begin{center}
\renewcommand{\arraystretch}{1.4}
\begin{tabular}{ll}
\toprule
\textbf{Symbol} & \textbf{Meaning} \\
\midrule
$k = 0,\ldots,K$ & Internal time index (frame number), $K+1$ frames total \\
$t \geq 0$ & Diffusion time (OU noise level) \\
$a_k \in \R$ & Clean frame at internal time $k$ \\
$x_k \in \R$ & Noisy observation of frame $k$ at diffusion time $t$ \\
$\ba, \bx \in \R^{K+1}$ & Full clean / noisy trajectory \\
$\alpha \in (-1,1)$ & AR(1) contraction coefficient \\
$\sigma_0^2$ & Variance of initial frame $a_0$ \\
$\sigma_\eta^2$ & AR(1) noise variance (variance of each $\eta_k$) \\
$\sigma_\infty^2 = \sigma_\eta^2/(1-\alpha^2)$ & Stationary variance \\
$\dt = 1 - e^{-2t}$ & OU noise variance at time $t$ \\
$\Sigma_0$ & $(K+1)\times(K+1)$ covariance of clean chain \\
$\Sigma_t = e^{-2t}\Sigma_0+\dt I$ & Covariance of noisy observations \\
$\bmu_a = (\mu_0,\alpha\mu_0,\ldots,\alpha^K\mu_0)^\top$ & Mean of clean chain \\
$\bmu_t = e^{-t}\bmu_a$ & Mean of noisy observations \\
$M(a_{k+1}\mid a_k)$ & AR(1) transition kernel \\
$P_0(\ba)$ & Joint distribution of clean trajectory \\
$P_t(\bx)$ & Joint distribution of noisy trajectory \\
$\bm{S}(\bx,t) = \nabla_{\bx}\log P_t(\bx)$ & Joint score (correct object) \\
$s(x_k,k,t) = \nabla_{x_k}\log p_t(x_k\mid k)$ & Marginal score (old, incorrect for propagator) \\
\bottomrule
\end{tabular}
\end{center}

% ══════════════════════════════════════════════════════════════════════════════
\section{Discussion: Alternative Noise Models for Sequences}
\label{sec:noise_discussion}
% ══════════════════════════════════════════════════════════════════════════════

\subsection{What does it mean to apply diffusion noise to a sequence?}

In our derivation we applied the OU process \emph{independently and
identically} to each frame:
\[
  x_k = a_k\,e^{-t} + \sqrt{\dt}\,z_k,
  \qquad z_k \overset{\mathrm{i.i.d.}}{\sim} \N(0,1),
  \qquad k = 0,\ldots,K.
\]
This choice is natural but is not the only possibility. The choice of
noise model fundamentally changes $P_t(\bx)$, the precision matrix
$\Sigma_t^{-1}$, and therefore the structure of the joint score.
Below we compare three alternatives systematically.

\subsection{Case 1: i.i.d.\ OU per frame (our setting)}

The noise covariance is $\text{Cov}(\bm{z}) = I_{K+1}$, so:
\[
  \Sigma_t = e^{-2t}\Sigma_0 + \dt\,I_{K+1}.
\]
The noise term $\dt I$ is \emph{diagonal} — it adds no coupling between
frames. All inter-frame structure in $\Sigma_t$ comes entirely from the
signal term $e^{-2t}\Sigma_0$, which decays to zero as $t\to\infty$.
This is what makes the score coupling vanish at large $t$.

The joint distribution $P_t(\bx)$ is Gaussian for all $t \geq 0$,
regardless of whether $P_0$ is Gaussian. This is because convolving
any distribution with a Gaussian (the OU noise) yields a distribution
that is at least as smooth, and for Gaussian $P_0$ it remains Gaussian
exactly.

\subsection{Case 2: correlated OU noise across frames}

Suppose instead the noise has a covariance structure $C \neq I$:
\[
  \bx = e^{-t}\ba + \sqrt{\dt}\,C^{1/2}\bz,
  \qquad \bz \sim \N(\bm{0}, I_{K+1}).
\]
Then $\text{Cov}(\sqrt{\dt}C^{1/2}\bz) = \dt C$, and:
\[
  \Sigma_t^{(C)} = e^{-2t}\Sigma_0 + \dt\,C.
\]
If $P_0$ is Gaussian, $P_t$ remains Gaussian — the closure property
holds. However, the precision matrix becomes
$(\Sigma_t^{(C)})^{-1}$, which has a different sparsity structure than
in Case~1. In particular, if $C$ is itself dense (e.g.\ correlated
noise between consecutive frames), the noise term \emph{adds}
off-diagonal entries, potentially strengthening rather than erasing
inter-frame coupling in the score. This would be physically meaningful
if the noise represents a correlated perturbation (e.g.\ camera shake
affecting all frames simultaneously).

\begin{remark}
A natural choice for a sequence is a \emph{temporally correlated} noise
with covariance $C_{jk} = \rho^{|j-k|}$ for some $\rho \in (0,1)$.
This adds a second tridiagonal structure to $\Sigma_t$, and the
precision $(\Sigma_t^{(C)})^{-1}$ remains banded but with potentially
wider bandwidth.
\end{remark}

\subsection{Case 3: a single shared noise applied to the full sequence}

At the opposite extreme, one could apply a single scalar noise
$\xi \sim \N(0,1)$ to the entire trajectory:
\[
  x_k = a_k\,e^{-t} + \sqrt{\dt}\,\xi \quad \text{for all } k.
\]
Then $\text{Cov}(\bx\mid\ba) = \dt\,\mathbf{1}\mathbf{1}^\top$ (a rank-1
matrix, where $\mathbf{1} = (1,\ldots,1)^\top$). This is a \emph{highly
degenerate} noise: all frames receive exactly the same random
perturbation. Even for Gaussian $P_0$, the resulting $P_t(\bx)$ is
Gaussian, but:
\[
  \Sigma_t^{(\text{shared})}
  = e^{-2t}\Sigma_0 + \dt\,\mathbf{1}\mathbf{1}^\top,
\]
which is rank-$(K+1) + 1$ in general but structurally very different from
Case~1. The noise now \emph{introduces} inter-frame correlation rather
than destroying it, so the score at large $t$ does not simplify to
$-\bx/\dt$ — instead it retains a non-trivial coupling through the
rank-1 noise term.

\subsection{What if \texorpdfstring{$P_0$}{P0} is not Gaussian?}

If $P_0(\ba)$ is not Gaussian (e.g.\ a mixture of Gaussians, or a
distribution with heavy tails), the answer depends on the noise model:

\begin{itemize}
  \item \textbf{i.i.d.\ OU (Case 1):} $P_t(\bx)$ is the convolution of
    $P_0$ with a Gaussian kernel. For small $t$ this preserves the
    non-Gaussian structure of $P_0$; for large $t$ the Gaussian kernel
    dominates and $P_t \approx \N(\bm{0},I)$. At intermediate $t$ the
    score $\bm{S}(\bx,t) = \nabla_{\bx}\log P_t(\bx)$ is \emph{nonlinear}
    in $\bx$ — the linear formula $-\Sigma_t^{-1}(\bx-\bmu_t)$ no
    longer applies.

  \item \textbf{Mixture of Gaussians:} if
    $P_0(\ba) = \sum_c \pi_c\,\N(\ba;\bmu_c,\Sigma_c)$, then
    $P_t(\bx) = \sum_c \pi_c\,\N(\bx;\,e^{-t}\bmu_c,\,e^{-2t}\Sigma_c
    + \dt I)$ — still a mixture of Gaussians. The score is a
    \emph{weighted average} of the per-component scores, with
    weights depending on the posterior component probabilities given $\bx$.
    This is nonlinear in $\bx$ in general.
\end{itemize}

\begin{resultbox}
\textbf{Key takeaway.} The Gaussian AR(1) model with i.i.d.\ OU noise is
the \emph{unique} setting in this family where: (i) $P_t$ is Gaussian
for all $t$, (ii) the joint score is linear in $\bx$, (iii) the score
has a closed-form precision structure, and (iv) the Kalman smoother
provides the exact posterior. Any departure from Gaussianity in $P_0$
breaks (i)--(iv).
\end{resultbox}

% ══════════════════════════════════════════════════════════════════════════════
\section{Extension to Multivariate Frames}
\label{sec:multivariate}
% ══════════════════════════════════════════════════════════════════════════════

\subsection{Setup}

Suppose each frame is now a $d$-dimensional vector:
\[
  \ba_k \in \R^d, \qquad k = 0, 1, \ldots, K.
\]
The full trajectory is $\ba = (\ba_0^\top, \ba_1^\top, \ldots,
\ba_K^\top)^\top \in \R^{d(K+1)}$. The AR(1) dynamics become:
\[
  \ba_{k+1} = A\,\ba_k + \bm{\eta}_k,
  \qquad
  \bm{\eta}_k \overset{\mathrm{i.i.d.}}{\sim} \N(\bm{0},\, Q),
\]
where $A \in \R^{d\times d}$ is the transition matrix (replacing the
scalar $\alpha$) and $Q \in \R^{d\times d}$ is the noise covariance
matrix (replacing $\sigma_\eta^2$).

The initial distribution is $\ba_0 \sim \N(\bmu_0, \Pi_0)$ for some
$\bmu_0\in\R^d$ and positive definite $\Pi_0\in\R^{d\times d}$.

\subsection{Does \texorpdfstring{$P_0(\ba)$}{P0(a)} remain Gaussian?}

\begin{proposition}
For multivariate $\ba_k \in \R^d$ with the AR(1) dynamics above, the
joint distribution $P_0(\ba) = P_0(\ba_0, \ldots, \ba_K)$ is a
$d(K+1)$-dimensional Gaussian.
\end{proposition}

\begin{proof}
The argument is identical to the scalar case. The Markov factorisation
holds:
\[
  P_0(\ba_0,\ldots,\ba_K)
  = p_0(\ba_0)\prod_{k=0}^{K-1} M(\ba_{k+1}\mid\ba_k),
\]
where $M(\ba_{k+1}\mid\ba_k) = \N(\ba_{k+1};\,A\ba_k,\,Q)$.
Each factor is Gaussian in the joint variable $(\ba_k,\ba_{k+1})$,
and the product of Gaussians with a shared linear structure is Gaussian.
\end{proof}

\subsection{The covariance matrix in the multivariate case}

The mean and covariance of $\ba_k$ generalise directly:
\[
  \bmu_k = A^k\,\bmu_0,
  \qquad
  \Pi_k = A^k\Pi_0(A^k)^\top + \sum_{j=0}^{k-1}A^j Q (A^j)^\top.
\]
The cross-covariance between frames $j \leq k$ is:
\[
  \Cov(\ba_j, \ba_k) = A^{k-j}\,\Pi_j \;\in\; \R^{d\times d}.
\]
The full joint covariance $\Sigma_0 \in \R^{d(K+1)\times d(K+1)}$
is now a \emph{block matrix} with $d\times d$ blocks:
\[
  (\Sigma_0)_{jk} = A^{|j-k|}\,\Pi_{\min(j,k)} \in \R^{d\times d}.
\]

\subsection{The precision matrix: block-tridiagonal structure}

\begin{proposition}[Block-tridiagonal precision]
The precision matrix $\Sigma_0^{-1}$ in the multivariate AR(1) case
is \textbf{block-tridiagonal}, with $d\times d$ blocks:
\begin{align*}
  (\Sigma_0^{-1})_{00} &= \Pi_0^{-1} + A^\top Q^{-1} A, \\
  (\Sigma_0^{-1})_{kk} &= Q^{-1} + A^\top Q^{-1} A,
    \quad k = 1,\ldots,K-1,\\
  (\Sigma_0^{-1})_{KK} &= Q^{-1}, \\
  (\Sigma_0^{-1})_{k,k+1} = [(\Sigma_0^{-1})_{k+1,k}]^\top
    &= -Q^{-1}A, \quad k = 0,\ldots,K-1.
\end{align*}
\end{proposition}

The derivation is the same as the scalar case: expand the exponent,
read off block coefficients, and use the symmetry of the full matrix.
The scalar entries $1/\sigma_0^2$, $(1+\alpha^2)/\sigma_\eta^2$,
$-\alpha/\sigma_\eta^2$ are simply replaced by the $d\times d$ matrices
$\Pi_0^{-1}$, $Q^{-1}+A^\top Q^{-1}A$, $-Q^{-1}A$.

\subsection{The OU diffusion and joint score in the multivariate case}

Applying i.i.d.\ OU noise independently to each frame component:
\[
  \bx_k = e^{-t}\ba_k + \sqrt{\dt}\,\bm{z}_k,
  \qquad
  \bm{z}_k \overset{\mathrm{i.i.d.}}{\sim} \N(\bm{0}, I_d).
\]
The result carries through exactly:
\[
  P_t(\bx) = \N\!\left(\bx;\; e^{-t}\bmu_a,\;
    e^{-2t}\Sigma_0 + \dt\, I_{d(K+1)}\right),
\]
and the joint score is:
\[
  \bm{S}(\bx,t) = -\Sigma_t^{-1}(\bx - e^{-t}\bmu_a)
  \;\in\; \R^{d(K+1)}.
\]
The $k$-th \emph{block} of the score (a $d$-vector) is:
\[
  \bm{S}_k(\bx,t)
  = -\sum_{j=0}^{K}(\Sigma_t^{-1})_{kj}\,(\bx_j - e^{-t}\bmu_j)
  \;\in\; \R^d.
\]
Everything remains Gaussian, the score remains linear, and the
precision matrix remains block-tridiagonal (approximately, for $t>0$).
The Kalman smoother generalises to the vector case without modification.

\begin{remark}[What changes]
The main additional complication is that $A$ is now a matrix, which
may not commute with $Q$. The off-diagonal blocks $-Q^{-1}A$ are
not generally symmetric, so the full precision matrix is symmetric
as a whole (as required) but its blocks are not individually symmetric.
For the special case $A = \alpha I_d$ and $Q = \sigma_\eta^2 I_d$
(isotropic scalar-times-identity dynamics), everything reduces to $d$
independent copies of the scalar AR(1) case and the structure is
exactly as before, just replicated $d$ times.
\end{remark}

% ══════════════════════════════════════════════════════════════════════════════
\section{The Kalman Filter and Smoother}
\label{sec:kalman}
% ══════════════════════════════════════════════════════════════════════════════

\subsection{Why the Kalman smoother matters here}

In the Mézard formula (Theorem~\ref{thm:mezard}):
\[
  S_k(\bx,t) = \frac{e^{-t}\langle a_k\rangle_{\ba\mid\bx} - x_k}{\dt},
\]
the central object is $\langle a_k\rangle_{\ba\mid\bx}$: the posterior
mean of the \emph{clean} frame $a_k$ given the entire \emph{noisy}
trajectory $\bx = (x_0,\ldots,x_K)$.

Computing this posterior mean is \emph{exactly} the problem solved
by the \textbf{Kalman smoother}. The Kalman smoother is therefore
not merely a useful tool — it \emph{is} the explicit propagator.

\subsection{The Kalman filter: forward pass}

The Kalman filter processes the noisy observations sequentially,
maintaining a Gaussian belief over the current clean frame.

\medskip
\noindent\textbf{Initialisation.} Prior belief at time $0$:
\[
  p(a_0) = \N(a_0;\;\mu_0,\;\sigma_0^2).
\]

\noindent\textbf{Prediction step.} Given the posterior at time $k$,
$p(a_k\mid x_0,\ldots,x_k) = \N(a_k;\;\hat{a}_{k|k},\;P_{k|k})$,
predict the next frame using the dynamics:
\[
  \hat{a}_{k+1|k} = \alpha\,\hat{a}_{k|k},
  \qquad
  P_{k+1|k} = \alpha^2\,P_{k|k} + \sigma_\eta^2.
\]

\noindent\textbf{Update step.} Upon observing $x_{k+1}$, compute the
\emph{Kalman gain}:
\[
  \mathcal{K}_{k+1}
  = \frac{e^{-t}\,P_{k+1|k}}{e^{-2t}\,P_{k+1|k} + \dt},
\]
and update the belief:
\[
  \hat{a}_{k+1|k+1}
  = \hat{a}_{k+1|k}
    + \mathcal{K}_{k+1}\!\left(x_{k+1} - e^{-t}\hat{a}_{k+1|k}\right),
  \qquad
  P_{k+1|k+1}
  = (1 - e^{-t}\mathcal{K}_{k+1})\,P_{k+1|k}.
\]

The innovation $x_{k+1} - e^{-t}\hat{a}_{k+1|k}$ is the surprise in the
new observation relative to the prediction. The Kalman gain
$\mathcal{K}_{k+1}$ weights how much to update: large when the
prediction uncertainty $P_{k+1|k}$ is large, small when the noise
$\dt$ is large.

\subsection{The Kalman smoother: backward pass}

The filter gives $\hat{a}_{k|k}$ — the best estimate using only
observations up to time $k$. The \emph{smoother} improves this by
running a backward pass that incorporates future observations.

\medskip
\noindent\textbf{Smoother initialisation.} At the last frame:
\[
  \hat{a}_{K|K}^s = \hat{a}_{K|K}, \qquad P_{K|K}^s = P_{K|K}.
\]

\noindent\textbf{Backward recursion.} For $k = K-1, K-2, \ldots, 0$:
\[
  L_k = P_{k|k}\,\alpha\,P_{k+1|k}^{-1}
  \quad\text{(smoother gain)},
\]
\[
  \hat{a}_{k|K}^s
  = \hat{a}_{k|k} + L_k\!\left(\hat{a}_{k+1|K}^s - \hat{a}_{k+1|k}\right),
  \qquad
  P_{k|K}^s
  = P_{k|k} + L_k\!\left(P_{k+1|K}^s - P_{k+1|k}\right)L_k^\top.
\]

The output $\hat{a}_{k|K}^s = \langle a_k\rangle_{\ba\mid\bx}$ is the
exact posterior mean we need for the Mézard formula.

\subsection{Connection to the joint score}

The Kalman smoother provides a \emph{recursive} way to evaluate
$\bm{S}(\bx,t)$ without ever explicitly forming or inverting the
$(K+1)\times(K+1)$ matrix $\Sigma_t$. Instead, it computes
$\langle a_k\rangle_{\ba\mid\bx}$ for all $k$ in $\mathcal{O}(K)$
operations (two linear passes), versus $\mathcal{O}(K^3)$ for direct
matrix inversion.

More importantly, the smoother \emph{is} the propagator. The backward
recursion step:
\[
  \hat{a}_{k|K}^s
  = \hat{a}_{k|k} + L_k\!\left(\hat{a}_{k+1|K}^s - \hat{a}_{k+1|k}\right)
\]
says: the smoothed estimate at frame $k$ is updated from the smoothed
estimate at frame $k+1$ via the smoother gain $L_k$. This is the
propagator $P_t$ mapping $S_{k+1} \to S_k$ in closed form.

\begin{resultbox}
\textbf{The Kalman smoother is the exact propagator.}
For the Gaussian AR(1) model, the hypothesised propagator
$S_k \approx P_t(S_{k-1})$ is not an approximation — it is exact,
and $P_t$ is the Kalman smoother gain operator. Computing the joint
score for the full $K+1$ frame trajectory costs $\mathcal{O}(K)$
via the smoother, not $\mathcal{O}(K^3)$ via matrix inversion.
\end{resultbox}

% ══════════════════════════════════════════════════════════════════════════════
\section{Suggested Visualisations}
\label{sec:plots}
% ══════════════════════════════════════════════════════════════════════════════

The following plots and animations are suggested to build intuition
for the results in this document. Each targets a specific conceptual
point.

\subsection{Static plots}

\begin{enumerate}

  \item \textbf{Covariance vs precision matrix (side by side, heat map).}
    For $K = 4$, plot $\Sigma_0$ and $\Sigma_0^{-1}$ as colour-coded
    grids with numerical values in each cell. Use distinct colours for
    diagonal, off-diagonal, and zero entries. Vary $\alpha$ with a
    slider to show how the dense $\Sigma_0$ compresses to diagonal while
    $\Sigma_0^{-1}$ remains tridiagonal. \textit{Key message: inversion
    extracts direct couplings.}

  \item \textbf{$\Sigma_t^{-1}$ entries vs diffusion time $t$.}
    Plot the diagonal entry $(\Sigma_t^{-1})_{kk}$ and the
    off-diagonal entry $|(\Sigma_t^{-1})_{k,k+1}|$ as functions of
    $t \in [0, 3]$. Mark the asymptotic values at $t=0$ (tridiagonal
    limit) and $t\to\infty$ (diagonal limit). \textit{Key message:
    inter-frame coupling in the score decays as $e^{-2t}$.}

  \item \textbf{Joint score $S_0(x_0,x_1,t)$ vs $x_0$ for fixed $x_1$,
    compared with marginal score.}
    For $K=1$: plot both $S_0^{\text{joint}}(x_0, x_1, t)$ (blue)
    and $s^{\text{marg}}(x_0, t) = -(x_0-\mu_{t,0})/T_{00}$ (red dashed)
    as functions of $x_0$. Move $x_1$ with a slider.
    \textit{Key message: the joint score shifts vertically with $x_1$;
    the marginal score is blind to it.}

  \item \textbf{The $K=2$ precision matrix at various $t$.}
    Four panels at $t = 0, 0.5, 1.0, 2.0$: show $\Sigma_t^{-1}$ as
    a $3\times 3$ grid. \textit{Key message: off-diagonal entries shrink
    and matrix approaches $I$ as $t$ grows.}

\end{enumerate}

\subsection{Animations}

\begin{enumerate}

  \item \textbf{OU forward process on a trajectory.}
    Animate $K+1$ frames of a clean AR(1) sample as $t$ increases
    from $0$ to $3$: each frame independently drifts toward $0$ and
    spreads, while the frames remain correlated (they all undergo
    the same shrinkage). Show the marginal distribution at each frame
    as a Gaussian envelope, and overlay sample paths.
    \textit{Key message: noise is added independently per frame,
    but the trajectory remains correlated throughout.}

  \item \textbf{Posterior mean (Kalman smoother) vs observation.}
    Fix a noisy trajectory $\bx$. As $t$ varies, animate how the
    Kalman-smoothed estimate $\langle\ba\rangle_{\ba\mid\bx}$ moves.
    At $t=0$: $\langle a_k\rangle = x_k e^{t} \approx x_k$ (perfect
    denoising). At large $t$: $\langle a_k\rangle \to \mu_{a,k}$
    (prior dominates, observation is useless).
    \textit{Key message: the score ``denoising residual'' shrinks as
    noise overwhelms the signal.}

  \item \textbf{Joint score vector field on $(x_0, x_1)$ plane, varying $t$.}
    For $K=1$: quiver plot of $(S_0, S_1)$ at each point of a grid.
    At $t \approx 0$: arrows are tilted off-axis (strong coupling).
    At large $t$: arrows point radially inward (score $\approx -\bx/\dt$,
    fully decoupled). Animate $t$ continuously.
    \textit{Key message: the propagator structure of the score fades
    continuously as $t$ grows.}

\end{enumerate}

% ══════════════════════════════════════════════════════════════════════════════
\end{document}
% ══════════════════════════════════════════════════════════════════════════════
